{\footnotesize
        \begin{tabularx}{\linewidth}{|p{0.25\columnwidth}|X|}
            \hline 
            \textbf{Notions et contenus} & \textbf{Capacités exigibles} \\
            \hline 
            \multicolumn{2}{|l|}{\textbf{5.1.3 Conduction électrique dans un conducteur ohmique}} \\
            \hline
                Loi d'Ohm locale. Conductivité électrique.  & 
                Établir l'expression de la conductivité électrique à l'aide d'un modèle microscopique, l'action de l'agitation thermique et des défauts du réseau étant décrite par une force de frottement fluide linéaire.
Discuter de l'influence de la fréquence sur la conductivité électrique.
Établir l'expression de la résistance d'une portion de conducteur filiforme. \\
            \hline 
             Effet Hall.   & 
            Interpréter qualitativement l'effet Hall dans une géométrie parallélépipédique. \\
        \hline 
        Effet thermique du courant électrique : loi de Joule locale. & 
        Exprimer la puissance volumique dissipée par effet Joule dans un conducteur ohmique. \\ 
        \hline
        \end{tabularx}
\vspace{0.5cm}
}
